% =====================================================================
%  MAKALE İSKELETİ — Türkçe akademik makale (XeLaTeX / LuaLaTeX)
%  Kişisel Zeka Kurulumu · sablonlar/akademisyen/latex/makale.tex
%
%  DERLEME (Katman 2 — Tectonic kuruluysa):
%      powershell -File araclar\belge\tex2pdf.ps1 -Girdi makale.tex
%  ya da doğrudan:
%      tectonic makale.tex
%
%  ⛔ pdflatex ile DERLEME. Türkçe karakterler için XeLaTeX/LuaLaTeX gerekir
%     (fontspec, sistem fontları). Tectonic zaten XeTeX tabanlıdır.
%
%  DERGİ ŞABLONU VARSA: derginin .cls dosyasını BOZMA. Bu dosyayı kullanma;
%  onun yerine gövdeni `md2tex.ps1 -Parca` ile üretip derginin şablonuna yapıştır.
% =====================================================================

\documentclass[12pt,a4paper]{article}

% ---- Türkçe + font (XeLaTeX/LuaLaTeX) -------------------------------
\usepackage{fontspec}
\setmainfont{Times New Roman}[
  % Sistemde Times New Roman yoksa şu satırı kullan (Linux/az fontlu makine):
  % ] \setmainfont{TeX Gyre Termes}[
  Ligatures=TeX
]
\usepackage{polyglossia}
\setmainlanguage{turkish}     % tireleme ve "Kaynakça/Şekil/Tablo" başlıkları Türkçe olur

% ---- Sayfa düzeni ---------------------------------------------------
\usepackage[a4paper,margin=2.5cm,left=3cm]{geometry}
\usepackage{setspace}
\onehalfspacing                % 1.5 satır aralığı
\usepackage{indentfirst}
\setlength{\parindent}{1.25cm} % ilk satır girintisi

% ---- Sık gereken paketler ------------------------------------------
\usepackage{graphicx}
\usepackage{booktabs}          % düzgün tablo çizgileri
\usepackage{longtable}
\usepackage{amsmath,amssymb}
\usepackage{hyperref}
\hypersetup{colorlinks=true, linkcolor=black, citecolor=black, urlcolor=blue}

% ---- Künye ----------------------------------------------------------
\title{\bfseries Çalışmanın Başlığı Buraya Gelir:\\ Alt Başlık Varsa İki Nokta Sonrası}
\author{Ad Soyad\thanks{Kurum, Bölüm. E-posta: ornek@ornek.edu.tr}}
\date{\today}

\begin{document}
\maketitle

% ---- Öz / Abstract --------------------------------------------------
\begin{abstract}
\noindent Öz buraya yazılır: çalışmanın amacı, yöntemi, bulguları ve sonucu 150--250
kelimeyle özetlenir. Öz tek paragraftır, kaynak içermez, kısaltma kullanmaz.

\vspace{0.5em}
\noindent\textbf{Anahtar Kelimeler:} birinci kelime, ikinci kelime, üçüncü kelime.
\end{abstract}

% ---- Gövde ----------------------------------------------------------
\section{Giriş}

Konu buradan başlar. Alanyazına atıf şöyle yapılır \cite{yilmaz2024}; birden çok
kaynak verilecekse \cite{yilmaz2024,ozturk2023} biçiminde yazılır.

Şekil ve tablolara \textbf{elle numara yazma} --- etiket ver, referansla:
Tablo~\ref{tbl:ornek} bunun örneğidir.

\section{Yöntem}

Araştırmanın deseni, örneklem, veri toplama araçları ve analiz yöntemi burada anlatılır.

\begin{table}[htbp]
  \centering
  \caption{Örnek tablo başlığı (tablo başlığı ÜSTTE olur).}
  \label{tbl:ornek}
  \begin{tabular}{lrr}
    \toprule
    Değişken & Ortalama & Standart Sapma \\
    \midrule
    Birinci ölçüm  & 12,4 & 2,1 \\
    İkinci ölçüm   & 15,8 & 3,4 \\
    Üçüncü ölçüm   & 11,2 & 1,9 \\
    \bottomrule
  \end{tabular}
\end{table}

\section{Bulgular}

Bulgular yorumsuz sunulur. Sayısal değerlerde Türkçe yazımda ondalık ayırıcı
\emph{virgül}dür (12,4), binlik ayırıcı noktadır (1.250).

\section{Tartışma ve Sonuç}

Bulgular alanyazınla karşılaştırılır, sınırlılıklar açıkça yazılır, sonraki
çalışmalara öneri verilir.

% ---- Kaynakça -------------------------------------------------------
% Bu şablon kaynakçayı BİLEREK elle taşır (thebibliography): tek geçişte,
% biber/bibtex olmadan derlenir. Tectonic'in standart kipi harici bibliyografya
% aracı çalıştırmaz — bu yüzden en güvenli yol budur.
%
% .bib dosyası kullanmak istersen (Overleaf ya da tam TeX kurulumu varsa):
%   1) Bu bloğu sil, preamble'a şunu ekle:
%        \usepackage[style=apa,backend=biber]{biblatex}
%        \addbibresource{kaynakca.bib}
%   2) Buraya \printbibliography yaz.
%   3) Derleme sırası: xelatex → biber → xelatex → xelatex
% Word/PDF çıktısında .bib kullanmak için LaTeX gerekmez:
%   md2docx.ps1 -Kaynakca kaynakca.bib -AtifStili apa.csl

\begin{thebibliography}{9}

\bibitem{yilmaz2024}
Yılmaz, A. (2024). \emph{Belge üretim hatlarında biçim tutarlılığı}. Örnek Yayınevi.

\bibitem{ozturk2023}
Öztürk, B., \& Şahin, C. (2023). Çok yazarlı bir çalışma başlığı.
\emph{Örnek Dergi, 12}(3), 45--67. https://doi.org/10.0000/ornek.2023.12345

\end{thebibliography}

\end{document}
